\section{近期研究（2026 年 8--11 月）}
\paragraph{Gaussian Sculpting：将网格抽取纳入表面训练循环}
2026.08~\cite{2608_10602}。先优化高斯再离线提取表面，容易出现用于重建的表面与用于渲染的原语并不一致。Gaussian Sculpting 从当前符号距离场作可微等值面抽取，在得到的网格上实例化高斯，并约束渲染原语保持在表面附近。双层训练让外层改善表面，内层拟合固定几何附近的高斯。

原文用物体场景的 Chamfer 距离及多视图渲染检查几何一致性。内层训练构成主要计算成本。作者说明双层优化的耗时和显存将方法限制在物体尺度，反射材料和细纹理仍可能削弱几何线索。

\paragraph{CoMVS-GS：让 PatchMatch 与高斯深度相互监督}
2026.08~\cite{2608_18413}。CoMVS-GS 用多视图立体匹配的稠密点及方向直接初始化压扁高斯，无需先做网格。优化中，高斯渲染给 PatchMatch 提供起始深度与法线，精化的匹配深度反过来约束高斯几何，最后以可见性证据和 Delaunay 图割提取表面，不经 TSDF 体融合。

论文评测具有激光参考的 DTU 以及由 GauU-Scene V2、MatrixCity 裁出的三组户外子场景。对象级与户外几何参照不同，户外结论限于子集而非完整城市。方法使用有位姿图像及多视图匹配，弱纹理和匹配不足区域仍需要评价其几何稳定性。

\paragraph{TopoSurfel：让训练中的代理网格反向约束高斯 surfel}
2026.08~\cite{2608_20687}。离散高斯 surfel 能渲染细节，但各原语局部拟合可能不能形成连续表面。TopoSurfel 在训练中动态提取无额外参数的代理网格，并把得到的结构反馈给高斯 surfel 的方向和密度演化，让表面不只在训练后作为一次性导出。

在 Mip-NeRF 360，补充材料同时展示网格、高斯法向和渲染；NeRF-Synthetic 的纹理图是在几何固定后另做光度优化，不能把纹理效果计作几何改进。可微等值面分辨率受 GPU 内存约束，镜面及透明材料仍是困难条件。

\paragraph{用法线和稠密深度先验约束高斯几何}
2026.09~\cite{2609_36969}。光度损失在反光或缺少纹理的区域很难唯一确定表面。该研究将法线先验用于约束高斯协方差的方向，再以稠密深度约束像素级结构。作者同时比较不同 SfM 管线，单独考察初始化点云的质量如何传递到高斯训练结果。

论文用 Tanks and Temples 的 Train 和 Horse 以及自采室内外场景评测。原文报告，较低的平均重投影误差与较高 PSNR 和 SSIM 对应，但与 LPIPS 的关系较弱，COLMAP 在部分停车场和街景数据还未能重建出点云。这限制了单一初始化和图像指标的解释范围。

\paragraph{分开检验亮度与色度如何参与高斯结构形成}
2026.10~\cite{2610_00749}。标准 RGB 训练同时确定高斯几何与外观。该研究分别用亮度、色度与 RGB 优化，再冻结空间结构、按同一外观拟合协议重新估计颜色，通过留出视图对比哪一种监督形成的空间支持更易重建。作者还改变球谐阶数以考察方向外观容量。

主几何比较在 11 个标准场景每条件运行四次，作者观察到亮度监督生成的结构总体更接近 RGB 条件，而仅色度形成的结构差距较大。初始化均来自用灰度特征的 COLMAP；外观重拟合还使用全部非测试图，包括原训练的验证视图。因此结论限定在该初始化、拟合和测量协议，不等同于直接取得独立几何真值。

\paragraph{Sparse-GS2Mesh：用虚拟立体视图互补约束二维和三维高斯}
2026.10~\cite{2610_04203}。稀疏视图下只拟合输入图容易让高斯产生不完整的表面。Sparse-GS2Mesh 用极线深度暖启动，自动选取适当的立体基线渲染新视图，再用立体匹配提供几何监督，并通过 2DGS 与 3DGS 的共同正则提高一致性。

原文在 DTU 以 Chamfer 距离分别评估低重叠和高重叠条件。在低重叠并采用极线初始化时，报告相对既有最佳 SolidGS 改善约 15\%；高重叠一组的表述为有竞争力。这个百分比对应表面距离和特定输入协议，不能解释为图像质量或各类数据上的普遍提升。

\paragraph{PCAsplat：在三维原语邻域直接约束表面几何}
2026.10~\cite{2610_11011}。渲染只要求像素颜色一致时，高斯中心未必落在真实表面。PCAsplat 对每个中心的局部原语邻域作主成分分析，用特征值约束邻域的法向厚度和平面覆盖，并用最小特征方向给高斯法线提供对齐目标。这些损失在场景空间计算，与渲染损失联合优化。

论文在 DTU 与 Tanks and Temples 用 Chamfer 距离和 F1 评价高斯中心对表面的逼近，并演示以优化后的中心直接进行实例分割与 Poisson 重建。邻域 PCA 隐含局部近似平面的假设，交叉或相邻的薄表面可能落入同一邻域，不能把渲染指标看作这种情况下表面正确的证明。

\paragraph{TopSDF：用顶点相关性和网格能量约束三角面提取}
2026.11~\cite{2611_18708}。从高斯场的四面体网格提取等值面时，有限网格和简单的独立顶点模型可能形成不规则三角形。TopSDF 把顶点的符号距离联系成高斯随机场，以学习到的相关性更灵活地确定边上的表面交点；再用网格能量和相应的高斯增密一起约束重建。

原文在无界场景分别测网格渲染和几何表面，另在 DTU 的 15 个扫描验证物体几何。三角形质量、表面误差和新视角 PSNR 对应不同目标。潜在协方差和跨顶点能量查询增加训练计算与内存，补充材料仍展示反光表面波纹、深凹处过平滑等失败。

\paragraph{RAPID：用偏振与稀疏锚点预初始化二维高斯表面}
2026.11~\cite{2611_24404}。偏振图像提供表面朝向线索，但法向单独不能确定深度。该方法将多视图偏振法向和稀疏几何锚积分为深度，融合后按曲率采样，初始化二维高斯的位置、尺度和朝向。训练继续按表面证据插入、沿主轴分裂、在覆盖约束下合并并裁除原语。

作者在 PANDORA 三个场景和 PISR 四个实拍场景上评价，PISR 提供结构光参考网格。文中报告约 4.3 分钟的优化运行时间，是同表 PhyGaP 对照的约五分之一，并在四个相关场景中的三个同时改善 Chamfer 距离与 F 分数。这不是采集至成品的全部时间，输入也要求多视图偏振、标定和有效几何锚点。
